\documentclass[usenames,dvipsnames]{beamer}

\mode<presentation> {
\usetheme{Montpellier}
}

\usepackage{graphicx} % Allows including images
\usepackage{booktabs} % Allows the use of \toprule, \midrule and \bottomrule in tables
\usepackage[frenchb]{babel}
\usepackage[T1]{fontenc}
\usepackage[utf8]{inputenc}
\setcounter{tocdepth}{2}
\graphicspath{{../Figures/}}
\newcounter{numexos}%Création d'un compteur qui s'appelle numexos
\setcounter{numexos}{0}%initialisation du compteur
\newcommand{\exercice}[1]{%Création d'une macro ayant un paramètre
\addtocounter{numexos}{1}%chaque fois que cette macro est appelée, elle ajoute 1 au compteur numexos
\textcolor{DarkOrchid}{Exercice\,\thenumexos\,:}\,#1%Met en rouge Exercice et la valeur du compteur appelée par \thenumeexos
}

\setbeamertemplate{title page}{
    \centering

    \huge\bfseries
    Fondamentaux théoriques du machine learning

\begin{figure}[!h]
\includegraphics[width=8cm]{knn}
\end{figure}

    }

\title[FTML]{title}


\date{} % Date, can be changed to a custom date

\makeatletter
\let\@@magyar@captionfix\relax
\makeatother

\begin{document}

\begin{frame}
\titlepage % Print the title page as the first slide
\end{frame}

\section*{Introduction} % (fold)
\label{sec:section_name}

\begin{frame}{Organization}

    Repository of the course: \url{https://github.com/nlehir/FTML}.
    \\

    This repo will contain all material related to the course.
    \begin{itemize}
        \item lecture slides
        \item practical sessions guidelines and solutions, exercises
        \item projects descriptions
	\item etc
    \end{itemize}
    
    Organisation of the course:
    \begin{itemize}
        \item $8\times 2$ hours lectures.
        \item $16\times 2$ hours practical sessions.
    \end{itemize}
\end{frame}

\begin{frame}{General information}

    \begin{itemize}
        \item Lectures: you can ask questions during the lectures, which can be interactive.
	\item Some textbooks are listed in references.pdf
	\item You can reach me by email / discord if there is a specific notion that you are not comfortable with (in which case we might spend more time on it).
	\item Before most practical sessions and lectures I will send some exercises to prepare and we will correct them during the session.
    \end{itemize}
\end{frame}


\begin{frame}{Textbooks}

Textbooks are better than slides !

\url{https://github.com/nlehir/FTML/blob/master/references.pdf} 

\end{frame}

\begin{frame}
\frametitle{Overview of lecture 1} % Table of contents slide, comment this block out to remove it
\tableofcontents % Throughout your presentation, if you choose to use \section{} and \subsection{} commands, these will automatically be printed on this slide as an overview of your presentation
\end{frame}

\section{Introduction to AI and ML}

\begin{frame}{Definitions}

    We will start by the a short general discussion on \textbf{{AI and ML}}.

    Quick remarks:
    \begin{itemize}
	\item Machine learning (ML) is \textbf{{easier}} to define than AI. Several
    definitions of AI exists, and the term can be used in different contexts. However the practical usage of the term "ML" has been shifting a bit for the last few years.
	\item These topics are evolving at a fast pace, it is hard to be up to date on all their aspects. 
    \end{itemize}


    
\end{frame}

\begin{frame}
\frametitle{Example 1: digit classification}
\begin{figure}[!h]
\centering
\includegraphics[width=8cm]{mnist.png}
\caption{MNIST database \cite{lecun-mnisthandwrittendigit-2010}}
\label{mnist}
\end{figure}
\end{frame}

\begin{frame}
\frametitle{Example 2: robots}
\begin{itemize}
	\item Boston Dynamics robots (video)
    \item \url{https://www.youtube.com/watch?v=LikxFZZO2sk}
    \item \url{https://www.youtube.com/watch?v=g0TaYhjpOfo}
    \item \url{https://www.youtube.com/watch?v=tF4DML7FIWk&t=57s} 
    \item \url{https://www.youtube.com/watch?v=EezdinoG4mk} 
\end{itemize}
\end{frame}

\begin{frame}
\frametitle{Example 3: go player}
\begin{figure}[!h]
\centering
\includegraphics[width=8cm]{go.png}
\caption{Go game, won by AlphaGo in 2017 \cite{SilverHuangEtAl16nature}}
\label{mnist}
\end{figure}
\end{frame}

\begin{frame}{Definition}
   \begin{itemize}
       \item Example definition of AI : "the theory and techniques aiming at emulating intelligence".
       \item Hence defining AI requires definining "intelligence".
   \end{itemize}
\end{frame}

\begin{frame}{Turing test}
    \begin{itemize}
	\item It is reasonable to say that the problem of AI was born at the same
        time as that of \textbf{{Computer Science}}. One of the founders of
            Computer Science is Alan Turing (1912-1954). 
        \item In 1950, with the article \textbf{{"Computing machinery and intelligence"}} \cite{Turing2009}, Turing introduces the \textbf{{Turing test}}. 
	\item One of the forms of a Turing Test is a game in which a computer tries to behave like a human by answering questions.
    \end{itemize}
\end{frame}


\begin{frame}
   \begin{figure}[htpb]
       \centering
       \includegraphics[width=0.8\linewidth]{turing}
       \label{fig:turing}
   \end{figure} 
\end{frame}

\begin{frame}{Turing test}
    \begin{itemize}
        \item Importantly, Turing stresses that the terms "think", and "machines" are far from being \textbf{{unambiguous}}. 
	\item Many (most ?) questions regarding our own "intelligence" and "consciousness" as humans remain open today.
	\item In that respect, it is tricky to give a indisputable definition of AI.
    \end{itemize}
\end{frame}


\begin{frame}{Moravec Paradox}
   \begin{itemize}
       \item It is harder to emulate or simulate simple sensorimotor capacities than level abstract reasoning.
       \item ie : the sensorimotor system of an ant is harder for us to emulate than beating the world champion at chess (Deep Blue, 1997).
   \end{itemize} 
\end{frame}

\begin{frame}{Conscious reasoning}
    Broadly speaking, the first phases of research in AI were more dedicated to
    emulating the \textbf{{conscious brain.}} 
    \begin{itemize}
        \item logic programming (Prolog)
        \item expert systems
    \end{itemize}
\end{frame}

\begin{frame}{History : AI Winter (1970's)}
   \begin{itemize}
       \item In 1974 fundings for AI reseach started being cut by goverments because of unsuccessful projects.
        \item there were actually several so-called "AI winters".
   \end{itemize}
\end{frame}

\begin{frame}{Famous AI : Deep Blue (1997)}
   \begin{itemize}
       \item Beat the world chess champion Gary Kasparov in 1997.
           \item 
               \begin{figure}[htpb]
                   \centering
                   \includegraphics[width=0.3\linewidth]{blue}
                   \label{fig:blue}
		   \caption{{Deep Blue}}
               \end{figure}
            \item Despite its name, Deep Blue is \textbf{{not related}} to Deep
                learning !
   \end{itemize} 
\end{frame}

\begin{frame}{Decision trees}
    One of the underlying mechanisms of Deep Blue were \textbf{{Decision
    Trees}} / minimax algorithm.
    \begin{itemize}
	\item Let $p$ be the \textbf{{branching factor}} of the tree. (Here, the average number of children at each node)
	\item Let $d$ be the \textbf{{depth}} of the tree. 
    \end{itemize}
    \exercice{\textbf{Computation time}}
	What is the order of magnitude of the number of nodes in the tree ?
    \textbf{Conclusion}: exploring the complete tree is not possible as soon as the depth is not very small.
\end{frame}

\begin{frame}{2000's and 2010's}
    \begin{itemize}
        \item Around the 2000's, there has been a shift in the focus of AI.
        \item Machine learning and data science have become prominent research and industrial topics.
        \item We could argue that emulating the \textbf{{unconscious brain}} is
            more of a focus today.
        \item An important contemporary and future trend in research might be
            to connect research on emulating the conscious brain and research
            emulating the unconscious brain.
    \end{itemize}
\end{frame}

\begin{frame}{Interesting slides on AI history}

    \url{https://home.mis.u-picardie.fr/~furst/IA.php} 

\end{frame}

\begin{frame}{Machine learning: proposition of definitions}
    \begin{itemize}
        \item (proposed definition of) \textbf{learning}: "Modification of a behavior, based on a life experiment"
        \item \textbf{{Machine learning}}: "artificial intelligence which can learn
    and model some phenomena without being explicitely programmed".
        \item In Machine learning, some parameters are learned in an
            \textbf{semi-automatic way}, based on a \textbf{dataset} in order to solve a problem or to optimize a
            solution.
    \end{itemize}
\end{frame}

\begin{frame}{Classical programming vs ML}
    \begin{itemize}
	\item Classical program: predict the total amount of money spent, based on the number of fruits bought and the price of each individual fruit (a summation is enough)
	\item ML program: predict which fruit a person buys, based on a database of customers and on some information about this person (e.g. buying log, age, profession).
    \end{itemize}
\end{frame}

\begin{frame}{Use cases of ML}
    ML is useful for problems:
\begin{itemize}
    \item that we cannot solve direcly thanks to an explicit representation (such as the amount of money spent in the previous first example).
    \item that we could solve explicitely, but the computationnal ressources would be too heavy (e.g. molecular dynamics), so we \textbf{approximate} a solution with ML.
\end{itemize}

Sometimes, ML solutions to a problem are effective, but the underlying mechanism is not perfectly understood.
\end{frame}

\begin{frame}{Denomination}
    \begin{itemize}
    	\item The name "machine learning" is rather deceitful !
	\item It is no more machine driven than any algorithm or coffee machine.
	\item The denomination "statistical learning" is used in some contexts.
    \end{itemize}
\end{frame}

\begin{frame}{Machine learning}
   Machine learning has received attention and funding because it has reached
    state-of-the-art efficiency on several problems such as:
    \begin{itemize}
        \item computer vision
        \item spam classification
        \item machine translation
        \item speech recognition
        \item self-driving cars
    \end{itemize}
    \textbf{{Deep learning}}  is involved in several of these solutions.
\end{frame}

\begin{frame}{Example 1: ImageNet}
    \begin{itemize}
        \item A database of images (more than $15$M), hand-annotated in order to
            indicate what objects are present in the image. More than $20000$
            categories of images.
        \item Contest: ImageNet Large Scale Visual Recognition Challenge.
        \item The best classification (top-5 error rate) went from $25\%$ in 2011 to
            $\simeq 15.3\%$ in 2012.
        \item The technology used was deep learning, exploiting GPUs (AlexNet).
    \end{itemize}

    \url{https://paperswithcode.com/sota/image-classification-on-imagenet} 
\end{frame}

\begin{frame}{Example 2: AlphaGo}
   \begin{itemize}
       \item In 2015: beats a professionnal player. In 2017: beats the world
           champion.
        \item Uses several technologies: among them \textbf{{Deep reinforcement
            learning.}} 
        \item Improvements: AlphaGo Zero, trained without a database of played
            games. In 2017, AlphaZero beats AlphaGo Zero after $3$ days of
            learning.
   \end{itemize} 
\end{frame}

\begin{frame}{Example 3: AlphaFold}
    \begin{itemize}
        \item Goal: to predict the spatial configuration of proteins, from
            their DNA sequence.
        \item Achieves a breakthgough performance on the CASP challenge:
            \begin{itemize}
                \item 2018: more than $50\%$ GDT (Global distance test),
                    whereas it was $\leq 40\%$ before then.
                \item 2020: $92,4\%$ GDT. At a $\geq 90\%$ score, the method is
                    considered competitive with experimental methods..
            \end{itemize}
        \item \url{https://alphafold.ebi.ac.uk/} 
        \item Also based on Deep learning.
    \end{itemize}
\end{frame}

\begin{frame}{Deep learning $\subset$ Machine Learning $\subset$ AI}

	\begin{itemize}
		\item Machine learning is not all about deep learning.
		\item In this course, we will not focus much on deep learning, but rather on more elementary / fundamental tools. But some parts of the course will be dedicated to neural networks.
		\item Deep learning is better suited for a dedicated course that allows experimentation with the network architectures.
		\item However, the notions that we will study are definitley useful to master in a deep learning project.
		\item Globally speaking, ML has a strong \textbf{experimental} component in industrial applications.
	\end{itemize}
\end{frame}


\begin{frame}{ML revolution}
\begin{itemize}
    \item Technical progress in the computing and storage capacities
    \item Increase in amount of \textbf{available data}. According to IBM, $10^{18}$
        bytes are created each day.
    \item Progress in algorithmic methods to analyze the data.
\end{itemize}
    
\end{frame}

\begin{frame}{ML revolution II}

    \begin{itemize}
	\item The speed of evolution has increased again with LLMs (first Chat GPT release: 2022)
	\item Spectrum of activities between maths and computer science: the jobs linked to AI and ML evolve. We could argue that the balance is more shifted towards data engineering / devops / mlops / computer science in general
	    than 20 years ago.
	\item In this course, we will not focus on LLMs, or large models in general, but rather on \textbf{elementary} aspects of ML (also: I am not an LLM expert). We could almost call the course "classical ML" ?
    \end{itemize}

\end{frame}

\section{Learning paradigms}
\frame{\tableofcontents[currentsection]}

\subsection{Supervised learning}%
\label{sub:supervised_learning}

\begin{frame}{Supervised learning}

    Predict an output from an input.
    \begin{itemize}
    	\item discrete output: classification
	\item continuous output: regression
    \end{itemize}

    We learn from a dataset $D_n$ of $n$ \textbf{labeled examples}; 
\begin{equation}
    D_n =\{(x_i,y_i), i\in[1, n]\}
\end{equation}
    (Do not hesitate to ask for a clarification when a notation is unclear.)
\end{frame}

\begin{frame}{Supervised Learning: example of regression}
        \begin{itemize}
            \item $x$: age and height of a person. 
                Here $x$, is a \textbf{{vector}} containing $2$ \textbf{{features}}.
            \item $y$:  record on a $100$ meters track
        \end{itemize}
\end{frame}

\begin{frame}{Regression: raw inputs}
       \begin{tabular}{cl}  
         \begin{tabular}{c}
           \includegraphics[width=5.2cm]{X}
           \end{tabular}
           & \begin{tabular}{l}
             \parbox{0.5\linewidth}{%  change the parbox width as appropiate
Input data.
    }
         \end{tabular}  \\
\end{tabular}
\end{frame}

\begin{frame}{Regression: raw ouptuts}
       \begin{tabular}{cl}  
         \begin{tabular}{c}
           \includegraphics[width=4.2cm]{time}
           \end{tabular}
           & \begin{tabular}{l}
             \parbox{0.5\linewidth}{%  change the parbox width as appropiate
    \exercice{}
    \addtocounter{numexos}{-1}

    Our objective is to build  a function that predicts the output (100 meters track record) as a function of the \textbf{input} (height and age of the person).

    What order of magnitude of precision do we expect to obtain with this example ?

    }
         \end{tabular}  \\
\end{tabular}
\end{frame}


\begin{frame}{Random variables}

    In a machine learning problem, there is (almost) always a statistical error associated with a prediction. Our goal is to minimize it !
    \\

    Datasets should be seen as \textbf{sampled from a random variable}. Hence, \textbf{any quantity derived from the dataset} should be analyzed as a random variable as well.

\end{frame}

\begin{frame}{Histograms}
    \begin{figure}[htpb]
    	\centering
    	\includegraphics[width=0.8\textwidth]{histogram_age}
    	\label{fig:histogram_age}
    \end{figure}

\end{frame}

\begin{frame}{Histograms}
    \begin{figure}[htpb]
    	\centering
    	\includegraphics[width=0.8\textwidth]{histogram_height}
    	\label{fig:histogram_age}
    \end{figure}
\end{frame}

\begin{frame}{Histograms}
    \begin{figure}[htpb]
    	\centering
    	\includegraphics[width=0.8\textwidth]{histogram_time}
    	\label{fig:histogram_age}
    \end{figure}
\end{frame}

\begin{frame}
	We assume that there is a relationship between the input variable $X$ and the output variable $Y$, given by a mapping $f$. $Y=f(X)$. Importantly, $f$ might contains some \textbf{randomness}.

	Equivalently, we assume the existence of a \textbf{joint distribution} $\rho$, such that $(X, Y)\sim \rho$.
	\begin{itemize}
	    \item \textbf{{Objective: }} find a good \textbf{estimation} $ \tilde{f}$, of $f$.
	    \item $ \tilde{f}$ maps an input to an output (most often deterministically). For the input $x_i$, we predict $ \tilde{f}(x_i)$.
	    \item In order to measure the quality of $\tilde{f}$, we use \textbf{loss functions}.
	\end{itemize}

\end{frame}


\begin{frame}{Notations}

    Usually, in this course
    \begin{itemize}
    	\item as uppercase letter will often denote a random variable, e.g: $X$
	\item a lowercase letter will be a sample $x_i$ or a specific value $x$ of the variable.
    \end{itemize}

    Do not hesitate to ask for a clarification when a notation is unclear.

\end{frame}

\begin{frame}
   \begin{itemize}
       \item The loss function should be a measure of the \textbf{discrepancy} between our prediction $\tilde{f}(x_i)$ and the \text{label} $y_i$.
	\item For an individual sample, \textbf{{a}} discrepancy is the
        squared loss
	    \begin{equation}
		(\tilde{f}(x_i)-y_i)^2
	    \end{equation}
	\item Taking into account the whole dataset, the mean squared error (MSE) \textbf{{loss function}} writes:
	    \begin{equation}
		MSE(\tilde{f})= \frac{1}{n}\sum^{n}_{i=1} (\tilde{f}(x_i)-y_i)^2
	    \end{equation}
	    \item Other loss functions are possible, like the absolute loss:
	    \begin{equation}
		AL(\tilde{f})= \frac{1}{n}\sum^{n}_{i=1} |\tilde{f}(x_i)-y_i|
	    \end{equation}
   \end{itemize} 
\end{frame}

\begin{frame}{Loss function}
   \begin{itemize}
       \item The loss function is a \textbf{{real number}} measuring the relevance of a \textbf{{collection}} of parameters ($ \tilde{f}$ is defined by these parameters.)
	\item The number of parameters depends on the situation, and varies between $1$ (e.g. for a simple linear model) and millions (e.g. for some deep neural networks).
   \end{itemize}

    \textbf{To what subset of functions does $ \tilde{f}$ belong ?}
\end{frame}


\begin{frame}{Optimization}
    We will often look for functions $\tilde{f}$ such that the loss is a small as possible (empirical risk minimization).
\begin{figure}[htpb]
    \centering
    \includegraphics[width=0.5\linewidth]{function}
    \label{fig:function}
\end{figure}    
\end{frame}


\subsection{Unsupervised learning} % (fold)
\label{sub:unsupervised_learning}

\begin{frame}
\frametitle{Unsupervised Learning}
\begin{itemize}
	\item From a number of samples ${x_i}$, we want to retrieve information on
        their structure: \textbf{{modelisation}}.
    \item
        \begin{itemize}
        \item density estimation
        \item clustering
        \item dimensionality reduction
    \end{itemize}
\end{itemize}
\end{frame}


\begin{frame}
\frametitle{Density estimation}
\begin{itemize}
    \item density estimation: anaylze the distribution of the data.
\begin{figure}[htpb]
    \centering
    \includegraphics[width=0.5\linewidth]{bimodal}
    \label{fig:bimodal}
\end{figure}

	\item Question : how do we constrain the distribution to fit the data ? (parametric or non parametric models)
	\item Loss functions are also used, although not as intuitive to define.
\end{itemize}
\end{frame}


\begin{frame}{Clustering}
\begin{figure}[htpb]
	\centering
	\includegraphics[width=0.7\textwidth]{it_3_assign_voronoi}
	\caption{K-means clustering}
	\label{fig:if_3_assign_voronoi}
\end{figure}
\end{frame}

\begin{frame}{Dimensionality reduction}
\begin{figure}[htpb]
	\centering
	\includegraphics[width=0.65\textwidth]{tsne}
	\caption{t-SNE projection of the MNIST dataset.}
	\label{fig:}
\end{figure}
\end{frame}

\subsection{Reinforcement learning} % (fold)
\label{sub:reinforcement_learning}

\begin{frame}
\frametitle{Reinforcement Learning}
\begin{itemize}
	\item A \textbf{more general paradigm} that describes an \textbf{agent} in a \textbf{world}.
	\item The standard formalization was the one proposed by Richard Sutton \cite{Sutton1998}.
	\item At each time, the world is in a state $s$. An agent perfoms an \textbf{action} $a$ according to a \textbf{policy} $\pi$. When performing an action, the agent receives an \textbf{reward} $r$.
\end{itemize}
\end{frame}

\begin{frame}{Reinforcement Learning applications}
    \begin{itemize}
	\item	Example : some Chessplayers, AlphaGo, automatic vacuum cleaner,
        self-driving cars (trajectory optimization, parking, etc.)

    \item    healthcare: \url{https://arxiv.org/pdf/1908.08796.pdf} 
        \item \url{https://deepmind.com/blog/article/alphago-zero-starting-scratch} 

    \end{itemize}
    
\end{frame}

\begin{frame}
\frametitle{Reinforcement Learning}
\begin{itemize}
	\item At each time, the world is in a state $s$. An agent perfoms an \textbf{action} $a$ according to a \textbf{policy} $\pi$. When performing an action, the agent receives an \textbf{reward} $r$.
	\item The agent wants to learn an \textbf{optimal policy}, meaning the policy that maximises its reward (or rather the \textbf{aggregation} of rewards).
\end{itemize}
\end{frame}

\begin{frame}
\frametitle{This paradigm has many variants}
\begin{itemize}
	\item 	State $s$, action $a$, policy $\pi$, reward $r$.
	\item	Is the policy \textbf{deterministic }? Is it \textbf{stochastic } ? 
	\item	Does the agent have a \textbf{model} of its environment ? 
	\item	How many steps ahead should the agent look ? 
	\item	All these conditions change the way the problem should be addressed and solved. The \textbf{Bellman equations} rule the computation of the optimal policy.
\end{itemize}
\end{frame}

\subsubsection{Exploration vs exploration} % (fold)
\label{ssub:exploration_vs_exploration}

\begin{frame}
\frametitle{Example problem}
\exercice{\textbf{Exploration and exploitation}}

	Should an agent explore its environment more or exploit what it has learnt so far ?
\end{frame}


\begin{frame}
\begin{itemize}
	\item These paradigms can be mixed. For instance,
	\begin{itemize}
		\item unsupervised learning can be used in a supervised learning problem (semi-supervised learning)
		\item unsupervised learning and supervised learning can be used in a reinforcement learning problem
		\item \textbf{self-supervised learning} creates labels from the dataset itsself. 
	\end{itemize}
	\item Given a problem and a dataset, formalizing a problem as a supervised learning or unsupervised learning task is a \textbf{choice}.
	\item Often, obtaining labels for all samples is a hard, long, or even impossible process: e.g. labelling millions of images. Importantly, the dataset collection, preparation or preprocessing is at least as important as the algorithms that are run on it !
\end{itemize}
\end{frame}

\subsection{Examples}%
\label{sub:examples}

\begin{frame}{Example I}

    Predict the winning team of an NBA game at half-time.
\begin{itemize}
    \item Dataset: $15$ years of games (comments, text): approximately $17000$
        games.
    \item The dataset is preprocessed to have as an input a time-series: each
        time contains the score \textbf{{and}}  $10$ technical features (rebounds, etc.).
        So for each time the dimension is $11$. Each game is a matrix of size $
        1440\times 11$, reorganized as a line vector.
    \item Output: Receiving team wins or looses (classification)
    \item Evaluation metric: classification error ("0-1" loss).
\end{itemize}
    
\end{frame}

\begin{frame}{Example II}
    Predict the quantity of oil in a rock.
    \begin{itemize}
        \item Input: tomographic image of a rock.
        \item Output: material of the rock, average presence of residual oil in
            the rock (regression).
    \end{itemize}
\end{frame}

\begin{frame}{Example III: what type of learning is it ?}

    \begin{itemize}
    	\item Input: database of clients of an insurance comrany
	\item Output: groups of clients, based on customer behavior, in order to propose targetted products for the identified groups.
    \end{itemize}

\end{frame}


\begin{frame}{Example IV}
    Detect issues in wind farms.
    \begin{itemize}
        \item Input: sensors on the wind turbine (wind direction, air
            temperature, electric tension, rotation speed, component
            temperature, etc.) as a time series. Each step represents $10$
            minutes (several years).
        \item Output: Power generated by the turbine (regression)
        \item Evaluation metric: MAE (mean absolute error).
    \end{itemize}
    
\end{frame}


\section{What makes ML a hard problem ?}
\frame{\tableofcontents[currentsection]}

\subsection{Overfitting}

\begin{frame}{Overfitting}
    Can we \textbf{generalize} the information learned from a dataset ?
\end{frame}

\subsection{Curse of dimensionality}%
\label{sub:curse_of_dimensionality}

\begin{frame}{Curse of dimensionality}
   These two numbers are important in machine learning:
    \begin{itemize}
        \item $n$: number of samples
        \item $d$: dimensionality (number of features) of a unique sample
    \end{itemize}

    Both can be large and prohibitive for some algorithms.
\end{frame}

\begin{frame}{Curse of dimensionality}
   Two numbers are important in machine learning:
    \begin{itemize}
        \item $n$: number of samples
        \item $d$: dimension (number of features) of a unique sample
    \end{itemize}

    Both can be large and prohibitive for some algorithms.

    \exercice{\textbf{Example:}} what is the algorithmic complexity of the
    inversion of an invertible matrix $A\in \mathbb{R}^{d, d}$ ?
\end{frame}

\begin{frame}{Curse of dimensionality}

    \begin{itemize}
        \item $n$ is large when the dataset has many samples.
        \item $d$ is large if each sample has many features:
            \begin{itemize}
                \item image
                \item DNA sequence
                \item text
                \item audio/video file
            \end{itemize}
    \end{itemize}
\end{frame}

\begin{frame}{Local averaging}

    The curse of dimensionality is the reason why \textbf{{local averaging}}
    does not work in large dimensions.
    
\end{frame}

\begin{frame}{Local averaging}
    \begin{figure}[htpb]
        \centering
        \includegraphics[width=0.8\linewidth]{knn10}
        \label{fig:name}
    \end{figure}
\end{frame}


\begin{frame}{Local averaging}
    \begin{figure}[htpb]
        \centering
        \includegraphics[width=0.8\linewidth]{knn100}
        \label{fig:name}
    \end{figure}
\end{frame}

\begin{frame}{Local averaging}
    \begin{figure}[htpb]
        \centering
        \includegraphics[width=0.8\linewidth]{knn10000}
        \label{fig:name}
    \end{figure}
\end{frame}


\begin{frame}{Local averaging}
    \begin{figure}[htpb]
        \centering
        \includegraphics[width=0.8\linewidth]{knn100000}
        \label{fig:name}
    \end{figure}
\end{frame}

\begin{frame}{In d dimensions}

    \begin{figure}[htpb]
        \centering
        \includegraphics[width=0.8\linewidth]{knnd}
        \label{fig:knnd}
    \end{figure}
    
\end{frame}



\subsection{Optimization}%
\label{sub:optimization}

    \begin{frame}{Optimization}

        We often have to optimize functions that are \textbf{{hard}} to
        optimize.

        \begin{itemize}
            \item non convexity
            \item non linearity
	    \item high dimension
            \item instability of estimators
        \end{itemize}
    \end{frame}

    \begin{frame}{Convex function}
	\begin{figure}[htpb]
		\centering
		\includegraphics[width=0.8\textwidth]{convex}
		\label{fig:convex}
	\end{figure}
    \end{frame}

    \begin{frame}{Non-convex function}
	\begin{figure}[htpb]
		\centering
		\includegraphics[width=0.8\textwidth]{nonconvex}
		\label{fig:convex}
	\end{figure}
    \end{frame}


\begin{frame}{The state of ML}
    ML is an active research area. One of the biggest effort is trying to understand the performance of deep learning.
    \begin{itemize}
        \item they are non convex
	\item but we experimentally manage to optimize them so that they have a state-of-the-art generalization error
    \end{itemize}
    In the case of neural networks with $1$ hidden layer, the convergence of gradient descent to a global minimizer of the objective function has been proven in 2018 in a specific case (large number of hidden neurons), using tools from optimal transport theory (\cite{chizat2018global}).
\end{frame}

\begin{frame}{The double descent phenomenon}
    \begin{figure}[htpb]
        \centering
        \includegraphics[width=1\textwidth]{doubledescent}
	\caption{[Belkin et al, 2019]}
        \label{fig:}
    \end{figure}
\end{frame}


\begin{frame}[allowframebreaks]
\frametitle{References}
\bibliographystyle{apalike}
\bibliography{epitech.bib} 
\end{frame}

\end{document} 
